\documentclass[twoside]{projektInzynierskiMS1z}
\usepackage{polski}
\usepackage[utf8]{inputenc}
\usepackage{amsmath}
\usepackage{graphicx}
\usepackage{url}
\usepackage{setspace}
\usepackage[colorlinks=true,linkcolor=black,citecolor=black,urlcolor=black]{hyperref}
%\drukJednostronny

%% naglowek bibliografii po polsku
\renewcommand{\refname}{Literatura}

%% tytul, promotor i autorzy
\title{Projekt i implementacja systemu symulacji oraz automatycznego
zarządzania ruchem kolejowym w środowisku webowym}

\promotor{dr inż. Rafał Brociek}
\rodzajProm{PROWADZĄCY PRACĘ}

\specjalnosc{Inżynieria Analizy Danych}
\ustawKierunekW{informatyka}
\jednostkaUczelni{WYDZIAŁ MATEMATYKI STOSOWANEJ}
\miasto{GLIWICE}
\rok{2026}

%% kazdy autor: imie nazwisko, nr albumu, opis wkladu
\autor{Martyna SZANDAŁA}{311093}{Opis wkładu}
\autor{Natalia TOMALA}{311109}{Opis wkładu}
\autor{Piotr DUSIŃSKI}{311489}{Opis wkładu}

%%%%%%%% zestaw komend opcjonalnych
%\NumeryNaPoczatku
%\subsectionWzory
%\sectionWzory
%\rozdzialy
%\literowaNumeracjaDodatkow
%\rzymskaNumeracjaDodatkow

\newtheorem{tw}{Twierdzenie}
\newtheorem{twa}{Twierdzenie}
\newtheorem{dd}{Definicja}

%%%%%%%%%%%%%%%%%%%%%%%%%%% strona druga -- streszczenia
\drugaStrona

\strDrNap{Streszczenie}

Przedmiotem niniejszej pracy jest zaprojektowanie oraz implementacja
wielowarstwowej aplikacji webowej służącej do interaktywnej symulacji
i monitorowania wirtualnego ruchu kolejowego. Głównym celem projektu było
stworzenie narzędzia, które w przystępny sposób obrazuje złożone procesy
zarządzania infrastrukturą kolejową, w tym działanie blokad liniowych oraz
procedury reagowania na sytuacje awaryjne.

\strDrNap{Słowa kluczowe}

symulacja ruchu kolejowego, aplikacja webowa, baza grafowa, algorytm A*,
Memgraph, rerouting

\strDrNap{Thesis title}

Design and implementation of a web-based railway traffic simulation
and automatic management system

\strDrNap{Abstract}

The subject of this thesis is the design and implementation of a multi-layered
web application for the interactive simulation and monitoring of virtual railway
traffic. The main goal of the project was to create a tool that, in an accessible
way, illustrates the complex processes of railway infrastructure management,
including the operation of line blocks and procedures for responding to emergency
situations.

\strDrNap{Keywords}

railway traffic simulation, web application, graph database, A* algorithm,
Memgraph, rerouting

\koniecDrugaStrona

\begin{document}

%% interlinia 1,5 dla tekstu podstawowego (zgodnie z formularzem Z4-PU12)
\onehalfspacing

%% UWAGA: naglowek "Wstep" (rozdzial 1) jest generowany automatycznie przez
%% klase zaraz po spisie tresci -- dlatego tu NIE wpisujemy \section{Wstep}.
%% Ponizszy tekst stanowi tresc wstepu.

Niniejszy projekt inżynierski koncentruje się na opracowaniu i implementacji
interaktywnej aplikacji webowej służącej do symulacji oraz monitorowania
wirtualnego ruchu kolejowego. Realizacja projektu obejmowała przygotowanie
koncepcji logiki biznesowej, zaprojektowanie interfejsu graficznego oraz
wdrożenie wielowarstwowej architektury systemu. Głównym celem aplikacji jest
stworzenie środowiska (przypominającego interaktywną makietę kolejową), które
w sposób przystępny i wizualny obrazuje skomplikowane procesy zarządzania ruchem
pociągów, takie jak działanie blokad liniowych, czy reagowanie na awarie
infrastruktury.

W pracy szczegółowo opisano cały proces realizacji projektu -- od analizy
wymagań funkcjonalnych, poprzez projektowanie architektury opartej na strukturach
grafowych i interfejsu, aż po implementację specyficznych algorytmów nawigacyjnych
i testowanie systemu. Ważnym elementem cyklu życia projektu było zastosowanie
zwinnych metodyk wytwarzania oprogramowania (ang. \emph{Agile}), co umożliwiło
elastyczne planowanie i etapowe rozszerzanie skali symulacji. Zrealizowane
rozwiązanie jest skierowane zarówno do pasjonatów kolejnictwa, jak i studentów
kierunków związanych z transportem, stanowiąc narzędzie ułatwiające zrozumienie
mechanizmów zapobiegających kolizjom i optymalizujących trasy przejazdów. Projekt
ten stanowi również solidną podstawę do dalszej rozbudowy o zewnętrzne źródła
danych.

\subsection{Geneza}

W ostatnich latach można zaobserwować silny trend cyfryzacji w sektorze transportu
publicznego, w tym wdrażanie koncepcji ,,cyfrowych bliźniaków'' (ang. \emph{digital
twins}) do symulowania zachowań rzeczywistej infrastruktury~\cite{digitaltwin}.
Wraz z tym zjawiskiem rośnie zapotrzebowanie na narzędzia edukacyjne i analityczne,
które w czytelny sposób tłumaczą zasady bezpieczeństwa na kolei. Inspiracją do
stworzenia aplikacji była obserwacja, że procesy decyzyjne dyspozytorów oraz
mechanizmy zabezpieczające tory są dla przeciętnego pasażera całkowicie niewidoczne,
a ich zrozumienie z samych podręczników jest trudne. Istniejące na rynku
i ogólnodostępne rozwiązania posiadają w tym zakresie pewne luki.

Popularne serwisy dla pasażerów (np. Portal Pasażera PKP PLK) skupiają się wyłącznie
na prezentacji aktualnych opóźnień i rozkładów jazdy, nie dając wglądu w przyczyny
tych zjawisk ani nie posiadając warstwy symulacyjnej. Z kolei zaawansowane symulatory
kolejowe (takie jak MaSzyna) kładą nacisk na trójwymiarowy realizm i perspektywę
maszynisty prowadzącego pojedynczy pojazd. Wymagają one wydajnego sprzętu i nie
pozwalają na spojrzenie na sieć kolejową z perspektywy makro -- jako na złożony
system logistyczny. Dostępne amatorskie symulatory urządzeń sterowania ruchem
(np. ISDR) są natomiast programami komputerowymi o wysokim progu wejścia, opartymi
na skomplikowanych, technicznych interfejsach odzwierciedlających prawdziwe pulpity
kostkowe. Brakowało na rynku nowoczesnego, przeglądarkowego narzędzia, które
łączyłoby rzeczywiste lub symulowane dane z przystępną symulacją na mapie
dwuwymiarowej.

\subsection{Odniesienie do rynku}

Rynek systemów informatycznych dla kolei (ang. \emph{Railway Management Systems})
dynamicznie rośnie, co jest napędzane potrzebą modernizacji infrastruktury
i zwiększania przepustowości linii. Prognozy rynkowe wskazują, że globalna wartość
tego sektora będzie w najbliższych latach systematycznie rosnąć, wymuszając
kształcenie nowych kadr w branży logistycznej~\cite{railwaymarket}.

Obecnie na rynku profesjonalnym dominują zaawansowane, zamknięte systemy
wykorzystywane przez zarządców infrastruktury (np. System Wspomagania Dyżurnego
Ruchu -- SWDR). Są to jednak narzędzia komercyjne, do których dostęp mają wyłącznie
uprawnieni pracownicy kolei. Z drugiej strony, w sektorze edukacyjnym i akademickim
brakuje lekkich, darmowych rozwiązań pozwalających na testowanie scenariuszy
kryzysowych na torach.

Większość dostępnych projektów akademickich skupia się na suchej analizie danych
algorytmicznych, pomijając aspekt wizualny. Projektowana aplikacja wypełnia tę niszę.
Zapewnia ona środowisko symulacyjne oparte na realnej siatce połączeń województwa
śląskiego, w którym użytkownik może obserwować interakcje między pociągami,
a administrator ma możliwość celowego wywoływania perturbacji w ruchu w celu badania
reakcji systemu.

\subsection{Cechy wyróżniające aplikację}

Projekt łączy intuicyjny interfejs z zaawansowaną architekturą logiki biznesowej.
Do głównych cech wyróżniających aplikację na tle prostych wizualizacji należą:

\begin{itemize}
\item \textbf{Zastosowanie grafowej bazy danych:} Wykorzystanie bazy działającej
w pamięci (Memgraph) pozwoliło na wysoce wydajne odwzorowanie sieci kolejowej.
Zapisanie stacji jako węzłów, a torów jako krawędzi, umożliwiło naturalne
i błyskawiczne operacje na strukturze sieciowej, niezbędne przy symulacjach
o dużej skali.

\item \textbf{Dynamiczny rerouting z algorytmem A*:} Aplikacja posiada wbudowany
silnik decyzyjny. W przypadku zablokowania odcinka toru, pociąg autonomicznie szuka
alternatywnej trasy. W tym celu zaimplementowano algorytm wyszukiwania najkrótszej
ścieżki A* (A-star), wykorzystując matematyczny wzór Haversine'a (obliczający
odległość na sferze) jako gwarantującą optymalność funkcję heurystyczną.

\item \textbf{Zwinne podejście środowiskowe:} W odróżnieniu od systemów zależnych
od zewnętrznych interfejsów API, projekt zrealizowano jako autonomiczną piaskownicę.
System opiera się na etapowo rozszerzanych danych statycznych 
% ( !Czy może być zmiana: pochodzących z Kolei Śląskich 
%  Stara wersja:
%  np. pobieranych
%  z OpenStreetMap
% )
pochodzących z Kolei Śląskich, co zapewnia niezawodność testów i uniezależnia działanie symulacji
od przerw w dostępie do sieci.

\item \textbf{Wbudowany silnik fizyki:} System nie tylko przesuwa znaczniki na mapie,
ale uwzględnia podstawowe zasady fizyki oraz logikę blokad liniowych, co weryfikuje
zajętość torów i zapobiega wirtualnym kolizjom. Mechanizm ten na bieżąco analizuje
potencjalne konflikty tras oraz bezwładność składów, co znacząco podnosi realizm
operacyjny. Dzięki takiemu rozwiązaniu aplikacja wiernie odzwierciedla zachowanie
pociągów, stanowiąc stabilne środowisko do przeprowadzania symulacji scenariuszy
awaryjnych.
\end{itemize}

\subsection{Wykorzystanie w praktyce}

Stworzona aplikacja może znaleźć szerokie zastosowanie w wielu realnych scenariuszach
edukacyjnych, demonstracyjnych i analitycznych, do których należą między innymi:

\begin{itemize}
\item \textbf{Prowadzenie zajęć z logistyki i transportu:} Wizualna demonstracja
w czasie rzeczywistym, jak działa system blokad liniowych i dlaczego pociągi muszą
zachowywać odpowiedni odstęp na szlaku.

\item \textbf{Analiza zarządzania kryzysowego typu what-if:} Dzięki funkcji
generowania awarii przez administratora, użytkownik zyskuje możliwość obserwacji,
jak zamknięcie jednego, kluczowego odcinka na Śląsku wpływa na kaskadowe opóźnienia
w całej sieci i jak reaguje algorytm wyznaczający objazdy.

\item \textbf{Wizualizacja priorytetów pociągów:} Narzędzie pozwala pokazać różnice
w obsłudze ruchu (np. w jakich sytuacjach pociąg towarowy musi zostać skierowany na
tor boczny, aby przepuścić skład zmierzający do stacji z jednoodcinkowego połączenia węzłów).

\item \textbf{Zrozumienie algorytmiki grafowej:} Aplikacja jest doskonałym narzędziem
dydaktycznym dla studentów informatyki, obrazującym praktyczne użycie bazy Memgraph
i algorytmu A* na rzeczywistych danych przestrzennych.
\end{itemize}

Dzięki swojej architekturze opartej na przejrzystym interfejsie graficznym aplikacja
obniża próg wejścia w złożony świat zarządzania ruchem kolejowym, pozwalając na obrazowe i
bezpieczne eksperymentowanie zarówno osobom nie związanym z koleją, jak i logistykom w uproszczony sposób.

\section{Założenia projektowe aplikacji}

Faza projektowa aplikacji poprzedzona została dogłębną analizą dostępnych rozwiązań
oraz zdefiniowaniem celów systemowych w oparciu o zasady inżynierii
wymagań~\cite{sommerville2016}. Ponieważ tworzenie zaawansowanego symulatora wymaga
precyzyjnego zbalansowania wydajności obliczeniowej z realizmem odwzorowania procesów,
zespół projektowy przyjął zwinne podejście do cyklu wytwarzania oprogramowania
(ang. \emph{Agile}). Metodyka ta pozwala na iteracyjne dostarczanie kolejnych
modułów funkcjonalnych oraz elastyczne reagowanie na problemy techniczne pojawiające
się w trakcie tworzenia projektu~\cite{agilemanifesto}.

Prace projektowe podzielono na etapy, zaczynając od definicji logiki biznesowej
i wymagań krytycznych, a kończąc na planowaniu modułów opcjonalnych. Taki podział
jest zgodny z dobrymi praktykami projektowania systemów informatycznych, gdzie
priorytetyzacja zadań pozwala na zachowanie spójności architektury przy rosnącej
skali projektu~\cite{martin2017clean}.

\subsection{Klasyfikacja wymagań systemowych}

Proces specyfikacji pozwolił na jasne wyodrębnienie rdzenia systemu (ang. \emph{core})
od funkcji rozszerzających. Zgodnie z teorią projektowania systemów zorientowanych na
dane, wymagania podzielono na kategorie uwzględniające ich wpływ na stabilność silnika
symulacyjnego:

\begin{itemize}
\item \textbf{Wymagania krytyczne:} Funkcjonalności niezbędne do realizacji
podstawowego modelu ruchu. Należą do nich: interaktywna wizualizacja węzłów i krawędzi
sieci, mechanizm zasilania danymi oraz system blokad liniowych. Bez tych elementów
system nie mógłby pełnić roli symulatora bezpieczeństwa ruchu kolejowego.

\item \textbf{Wymagania rozszerzone:} Moduły podnoszące wartość analityczną projektu,
takie jak algorytm dynamicznego wyznaczania objazdów (rerouting) oraz moduł
administratora do generowania zdarzeń losowych. Ich implementacja pozwala na
symulowanie rzeczywistych problemów występujących w sieciach transportowych,
wzbogacając znacząco wartość dydaktyczną całego projektu~\cite{sommerville2016}.

\item \textbf{Wymagania opcjonalne:} Funkcje o niższym priorytecie, których realizacja
jest uzależniona od wydajności końcowej wersji systemu (np. integracja z zewnętrznymi
interfejsami czy import rozbudowanych rozkładów jazdy). Ograniczenia projektu obejmują
świadome wykluczenia (np. rezygnacja z renderowania mapy całego kraju), mające na celu
utrzymanie wysokiej responsywności aplikacji webowej wyłącznie w środowisku przeglądarkowym.
\end{itemize}

\subsection{Opis kluczowych funkcjonalności systemu}

Na podstawie przeprowadzonej analizy wyodrębniono moduły funkcjonalne stanowiące trzon
aplikacji. Zespół zdecydował się na architekturę opartą na niezależnym silniku
symulacyjnym, co wymagało zaprojektowania następujących mechanizmów:

\subsubsection{Silnik symulacyjny i zabezpieczenia ruchu}

Sercem systemu jest moduł przeliczający pozycje pociągów na podstawie wbudowanego
modelu fizycznego. Kluczowym elementem jest integracja z systemem blokad liniowych,
który operuje na stanach krawędzi grafu. Rozwiązanie to znajduje oparcie w zasadach
działania urządzeń sterowania ruchem kolejowym (SRK), gdzie nadrzędnym celem jest
eliminacja ryzyka kolizji poprzez ścisłą kontrolę zajętości szlaków i odpowiednie
zależności odstępowe~\cite{srk}.

\subsubsection{Dynamiczny rerouting i algorytmy ścieżkowe}

Zidentyfikowano, że statyczne przypisanie tras jest niewystarczające w warunkach
dynamicznych awarii. Zastosowano więc mechanizm automatycznego wyznaczania dróg
alternatywnych oparty na algorytmie heurystycznym A* (A-star). Wykorzystuje on
przeszukiwanie przestrzeni stanów grafu, co w literaturze przedmiotu wskazywane jest
jako najskuteczniejsze narzędzie do optymalizacji przepływów w sieciach o strukturze
rozproszonej~\cite{astar1968}. Aby zagwarantować matematyczną poprawność i optymalność
wyznaczanych tras, algorytm zintegrowano ze wzorem Haversine'a, pełniącym funkcję
dopuszczalnej heurystyki szacującej odległość sferyczną między
stacjami~\cite{haversine}.

\subsubsection{Moduł administracyjny i analiza ,,what-if''}

Aby aplikacja miała charakter narzędzia badawczego, wprowadzono funkcję celowego
generowania awarii. Pozwala to na prowadzenie testów scenariuszowych (ang.
\emph{what-if}), co jest standardową praktyką w projektowaniu systemów krytycznych,
pozwalającą na identyfikację tak zwanych wąskich gardeł infrastruktury jeszcze przed
ich wystąpieniem w rzeczywistości~\cite{sommerville2016}.

\section{Instrukcja obsługi programu}
bla bla jadi jadi jada
% Pierwsze co widzi użyktownik
\subsection{Strona prezentacyjna} 
% Faktyczny panel symulacji
\subsection{Panel symulacyjno-projektowa}
% Panel do zarządzania użykownikami oraz rolami
\subsection{Panel administracyjna}

\section{Opis architektury programu}
W niniejszym rozdziale szczegółowo opisano stos technologiczny oraz założenia wzorców projektowych wykorzystanych w aplikacji. 
Architektura systemu opiera się na trzech głównych filarach:
\begin{itemize}
    \item Warstwa prezentacyjna (w mowie techniczno-korporacyjnej tzw. \verb|Frontend|) -- elementy interaktywne, widziane przez użytkownika po wejściu przez niego na adres strony aplikacji. Odbiorca za pomocą używania interfejsu dokonuje pożądanych przez niego zmian, przewidzianych w ramach funkcjonalności programu.
    \item Warstwa komunikacyjno-zarządzająca (ang. \emph{Backend} - są to procesy niezależne i niewidoczne dla użytkownika końcowego. Odpowiada za całą logikę aplikacji, komunikację z pozostałymi warstwami, niezbędne obliczenia oraz procesy uwierzytelniania i autoryzacji.
    \item Warstwa bazodanowa -- wszelkie szczegółowe czy też zmienne informacje są przechowywane w bazie danych. W warstwie tej są zapisane dane logowania, poszczególne stacje oraz elementy, które  są wymagane do poprawnego działania symulacji.
\end{itemize}
Do wprowadzenia obrazowego za pomocą filarów, należy dodać fundament, na którym stoją -- jest to środowisko wirtualizacyjne.  Wirtualizacja odpowiada za poszczególne połączenia między warstwami oraz ułatwienia pracy zespołowej. Powyższy wstęp do rozdziału ma na celu przedstawienie w skrócie, jakie zadania każda z warstw ma. Podrozdziały mają na celu wyjaśnienie w jak najbardziej zrozumiały i techniczny sposób działania komponentów składających się na całość opisywanego oprogramowania.
% Opisanie dockera i jego sieci-podsieci
\subsection{Architektura i środowisko}
% Opsianie architektury frontendu
\subsection{Warstwa prezentacyjna}
% Backend oraz to jak cały projekt komunikuję się z sobą
\subsection{Warstwa komunikacyjno-zarządzająca}
% Opisanie architektury bazy danych
% Źródła czemu ten rodzaj bazy lepszy, czemu ją wybraliśmy, jak powstała baza, źródła skąd powstawała, 
\subsection{Warstwa bazodanowa}


\section{Stos technologiczny}
W niniejszym rozdziale szczegółowo opisano narzędzia oraz technologie wybrane do
realizacji projektu interaktywnego symulatora ruchu kolejowego. Dobór stosu
technologicznego był podyktowany potrzebą zapewnienia wysokiej wydajności obliczeniowej
systemu, szczególnie w zakresie przetwarzania danych grafowych, oraz koniecznością
stworzenia responsywnego i intuicyjnego interfejsu użytkownika. Architektura systemu
została podzielona na niezależne warstwy: kliencką, serwerową oraz bazodanową, co
pozwala na zachowanie czystości kodu i łatwą rozbudowę poszczególnych modułów
w przyszłości.

\subsection{Warstwa kliencka}

Warstwa kliencka (ang. \emph{frontend}) odpowiada za wizualizację siatki połączeń
kolejowych na mapie geograficznej oraz prezentację stanów pociągów w czasie
rzeczywistym. Kluczowym wyzwaniem w tej warstwie było zapewnienie płynności animacji
przy dużej liczbie przemieszczających się obiektów, przy jednoczesnym utrzymaniu
niskiego zapotrzebowania na zasoby sprzętowe użytkownika.

\subsubsection{Svelte i SvelteKit}

Głównym narzędziem wybranym do budowy interfejsu użytkownika jest kompilator Svelte.
W odróżnieniu od tradycyjnych bibliotek, Svelte przenosi większość pracy na etap
budowania aplikacji, generując wysoce zoptymalizowany kod w czystym języku
JavaScript~\cite{svelte}. Pozwala to na uniknięcie narzutu obliczeniowego związanego
z ciągłym przeliczaniem wirtualnego modelu obiektu dokumentu (ang. \emph{Virtual
Document Object Model -- VDOM}). W przypadku symulatora, w którym pozycje wielu
pociągów muszą być stale aktualizowane z ułamkową precyzją, brak tego narzutu
drastycznie poprawia płynność działania aplikacji. Jako nadrzędny szkielet
aplikacyjny (ang. \emph{framework}) zastosowano SvelteKit, co ułatwiło zarządzanie
architekturą projektu oraz trasowaniem (ang. \emph{routing}).

\subsubsection{Leaflet}

Do renderowania interaktywnej mapy dwuwymiarowej wykorzystano bibliotekę Leaflet.
Charakteryzuje się ona wysoką wydajnością oraz stosunkowo niską wagą, co jest istotne
w kontekście aplikacji przeglądarkowych~\cite{leaflet}. Leaflet pozwala na
bezproblemową integrację z zewnętrznymi dostawcami kafelków mapowych (ang. \emph{tile
maps}), takimi jak OpenStreetMap~\cite{openstreetmap}, co umożliwiło osadzenie wirtualnej
sieci torów w rzeczywistym kontekście geograficznym. Biblioteka ta posiada rozbudowany
interfejs programistyczny dostosowany do dynamicznego manipulowania znacznikami
(reprezentującymi pociągi) oraz rysowania wektorowych linii szlaków, co stanowi
fundament wizualny całego symulatora.

\subsubsection{TypeScript}

Całość logiki po stronie klienta została napisana w języku TypeScript, który wprowadza
do JavaScriptu silne typowanie statyczne~\cite{typescript}. Ścisła integracja
TypeScriptu ze środowiskiem Svelte pozwala na rygorystyczne definiowanie struktur
danych przesyłanych z serwera (np. węzłów grafu, parametrów fizycznych składów).
Wczesne wykrywanie błędów na etapie kompilacji znacząco zwiększa stabilność kodu
i ułatwia zarządzanie rosnącą złożonością systemu informatycznego.

\subsubsection{Tailwind CSS}

Do budowy warstwy wizualnej interfejsu administracyjnego, paneli bocznych oraz
statystyk wykorzystano narzędzie Tailwind CSS~\cite{tailwind}. Jest to biblioteka
oparta na klasach użytkowych (ang. \emph{utility-first}), która umożliwia szybkie
i responsywne ostylowanie komponentów bezpośrednio w kodzie znaczników. Zastosowanie
tego rozwiązania wyeliminowało problem utrzymywania rozbudowanych i trudnych
w zarządzaniu plików ze stylami, przyspieszając tym samym prace nad oknami dialogowymi
używanymi do generowania awarii w systemie.

\subsubsection{Licencje}

Wszystkie użyte technologie frontendowe są udostępniane na darmowych licencjach
o otwartym kodzie źródłowym (ang. \emph{open-source}), z czego większość bazuje na
permisywnej licencji MIT. Pozwala to na ich pełne i legalne wykorzystanie w projektach
akademickich bez obaw o naruszenie własności intelektualnej.

\subsection{Warstwa serwerowa}

Warstwa serwerowa (ang. \emph{backend}) stanowi główne centrum obliczeniowe i decyzyjne
symulatora, realizujące zaawansowaną logikę biznesową, w tym wyznaczanie tras oraz
symulację fizyki ruchu pociągów. Architektura tej warstwy musiała zostać zaprojektowana
z myślą o wysokiej dostępności oraz zdolności do równoległego przetwarzania dziesiątek
operacji na sekundę. Z tego względu zrezygnowano z tradycyjnego podejścia
synchronicznego na rzecz rozwiązań wysoce zoptymalizowanych pod kątem operacji
wejścia-wyjścia, co gwarantuje płynność przepływu informacji w czasie rzeczywistym.

\subsubsection{Python i FastAPI}

Jako główny język programowania po stronie serwera wybrano Pythona, co było podyktowane
jego wszechstronnością oraz pozycją rynkowego standardu w środowiskach inżynierskich.
Bogaty ekosystem sprawdzonych bibliotek matematycznych oraz grafowych pozwala na
stabilną implementację złożonych algorytmów nawigacyjnych bez konieczności budowania
ich od podstaw. W roli nadrzędnego szkieletu programistycznego (ang. \emph{framework})
wykorzystano FastAPI, czyli jedno z najszybszych i najnowocześniejszych rozwiązań
dostępnych obecnie na rynku. Narzędzie to natywnie obsługuje asynchroniczność oraz
automatycznie generuje interaktywną dokumentację interfejsu programowania aplikacji
(ang. \emph{Application Programming Interface -- API}) w standardzie OpenAPI, co znacząco
usprawniło testowanie komunikacji pomiędzy serwerem a aplikacją kliencką~\cite{fastapi}.

\subsubsection{Warstwa przetwarzania asynchronicznego}

Z uwagi na specyfikę projektu, w którym parametry fizyczne i współrzędne
przemieszczających się pociągów muszą być aktualizowane w stałych interwałach
czasowych, konieczne było wdrożenie wydajnej warstwy przetwarzania asynchronicznego.
System wykorzystuje wbudowaną w język Python pętlę zdarzeń (ang. \emph{event loop}),
która umożliwia współbieżne przeliczanie danych dla wielu składów kolejowych bez
blokowania głównego wątku serwera~\cite{pythonasyncio}. Dzięki takiemu mechanizmowi
aplikacja może jednocześnie odpowiadać na nagłe zapytania użytkowników (np. żądanie
zablokowania toru) oraz bez przerw przesyłać zaktualizowane pozycje geograficzne do
warstwy interfejsu wizualnego.

\subsubsection{Bazy danych}

Warstwa trwałości danych odpowiada za przechowywanie oraz błyskawiczne udostępnianie
kluczowych informacji o strukturze sieci kolejowej. Z uwagi na fakt, że mapa połączeń
transportowych jest w swojej ostatecznej formie złożonym grafem, tradycyjne relacyjne
bazy danych okazały się nieoptymalne ze względu na wysoki koszt obliczeniowy zapytań
o wielokrotnie zagnieżdżone relacje~\cite{graphdatabases}.

\subsubsection{Memgraph}

Fundamentem systemu zarządzania danymi w projekcie jest grafowa baza danych
Memgraph~\cite{memgraph}. Przechowuje ona informacje w postaci wierzchołków
reprezentujących stacje kolejowe oraz krawędzi odzwierciedlających łączące je szlaki,
co stanowi bardzo naturalne i intuicyjne odwzorowanie fizycznej infrastruktury torowej.
Kluczową przewagą bazy Memgraph jest jej architektura oparta na przetwarzaniu danych
w całości w pamięci operacyjnej komputera (ang. \emph{in-memory}). Cecha ta pozwala na
wykonywanie złożonych zapytań przestrzennych i algorytmów przeszukiwania ścieżek
w czasie liczonym w pojedynczych milisekundach. Dodatkowo narzędzie to wspiera
ustandaryzowany język zapytań Cypher, który charakteryzuje się wysoką czytelnością
i ułatwia szybkie modelowanie warunków brzegowych niezbędnych do sprawnego działania
algorytmu wyznaczającego objazdy~\cite{cypher}.

\subsection{Konteneryzacja}

W celu zapewnienia wysokiej powtarzalności środowiska uruchomieniowego oraz ułatwienia
procesu wdrożenia (ang. \emph{deployment}) aplikacji na różnych infrastrukturach
serwerowych, w projekcie zdecydowano się na zastosowanie technologii konteneryzacji.
Podejście to pozwala na hermetyczne odizolowanie poszczególnych warstw systemu --
klienckiej, serwerowej oraz bazodanowej -- od bezpośredniej konfiguracji systemu
operacyjnego maszyny hosta~\cite{docker}. Dzięki takiemu rozwiązaniu wyeliminowano
powszechny w inżynierii oprogramowania problem niezgodności wersji bibliotek systemowych
oraz ułatwiono skalowanie całego rozwiązania. Konteneryzacja sprzyja również utrzymaniu
czystości w środowisku deweloperskim, gdyż wszystkie niezbędne zależności są instalowane
wyłącznie wewnątrz odizolowanych jednostek, co gwarantuje spójność działania aplikacji
niezależnie od platformy, na której zostanie uruchomiona.

\subsubsection{Docker}

Podstawowym narzędziem wykorzystanym do wirtualizacji na poziomie systemu operacyjnego
jest Docker. Pozwala on na zamknięcie każdej części aplikacji w niezależnym, lekkim
środowisku uruchomieniowym zwanym kontenerem (ang. \emph{container}). W kontekście
symulatora kolejowego przygotowano dedykowane obrazy (ang. \emph{images}) dla aplikacji
serwerowej napisanej w języku Python, warstwy klienckiej opartej na technologii Svelte
oraz bazy danych Memgraph. Każdy z tych obrazów definiuje kompletny zestaw bibliotek
i narzędzi niezbędnych do poprawnej pracy danego modułu. Wykorzystanie Dockera pozwala
na szybkie odtworzenie całego stosu technologicznego na dowolnym urządzeniu, co jest
kluczowe podczas pracy zespołowej oraz przy późniejszych testach integracyjnych systemu.
Dodatkowo, mechanizm warstwowości obrazów Dockerowych optymalizuje czas budowania
aplikacji, co znacząco przyspiesza proces iteracyjnego wprowadzania poprawek w logice
symulacji fizyki ruchu czy algorytmach dynamicznego wyznaczania objazdów.

\subsubsection{Docker Compose}

Do zarządzania architekturą składającą się z wielu współpracujących ze sobą kontenerów
wykorzystano narzędzie Docker Compose. Umożliwia ono zdefiniowanie całej infrastruktury
projektowej w postaci pojedynczego, deklaratywnego pliku konfiguracyjnego, co stanowi
realizację nowoczesnego paradygmatu infrastruktury jako kodu (ang. \emph{Infrastructure
as Code -- IaC})~\cite{infrastructureascode}. W pliku tym precyzyjnie określono
zależności pomiędzy poszczególnymi usługami, mapowanie portów sieciowych oraz
współdzielone woluminy danych (ang. \emph{volumes}), które zapewniają trwałość
informacji przechowywanych w bazie Memgraph nawet po zatrzymaniu kontenera. Docker
Compose tworzy dedykowaną, izolowaną sieć wewnętrzną, w której aplikacja serwerowa może
bezpiecznie komunikować się z bazą danych bez wystawiania jej portów na bezpośrednie
działanie czynników zewnętrznych. Dzięki temu uruchomienie całego, złożonego systemu
symulacyjnego sprowadza się do wykonania pojedynczego polecenia w konsoli, co
drastycznie obniża próg wejścia dla nowych osób pracujących nad projektem
i minimalizuje ryzyko błędów konfiguracyjnych.

%%%%%%%%%%%%%%%%%%%%%%%%%%% bibliografia -- recznie, wg Wzoru A (formularz Z4-PU12)
%% UWAGA: daty dostepu do stron WWW oraz dane pozycji nalezy zweryfikowac.
\begin{thebibliography}{20}

\bibitem{leaflet} Agafonkin V., Leaflet -- a JavaScript library for interactive maps, \url{https://leafletjs.com/} (dostęp: 30.06.2026).

\bibitem{agilemanifesto} Beck K. i in., Manifesto for Agile Software Development, 2001, \url{https://agilemanifesto.org/} (dostęp: 30.06.2026).

\bibitem{srk} Dąbrowa-Bajon M., Podstawy sterowania ruchem kolejowym. Funkcje, wymagania, zarys techniki, Oficyna Wydawnicza Politechniki Warszawskiej, Warszawa 2014.

\bibitem{docker} Docker Inc., Docker Documentation, \url{https://docs.docker.com/} (dostęp: 30.06.2026).

\bibitem{digitaltwin} Errandonea I., Beltrán S., Arrizabalaga S., Digital Twin for maintenance: A literature review, Computers in Industry, vol. 123, 2020, str. 103316.

\bibitem{astar1968} Hart P.E., Nilsson N.J., Raphael B., A Formal Basis for the Heuristic Determination of Minimum Cost Paths, IEEE Transactions on Systems Science and Cybernetics, vol. 4, nr 2, 1968, str. 100--107.

\bibitem{martin2017clean} Martin R.C., Clean Architecture: A Craftsman's Guide to Software Structure and Design, Prentice Hall, Boston 2017.

\bibitem{memgraph} Memgraph, Memgraph Documentation, \url{https://memgraph.com/docs} (dostęp: 30.06.2026).

\bibitem{typescript} Microsoft, TypeScript Documentation, \url{https://www.typescriptlang.org/docs/} (dostęp: 30.06.2026).

\bibitem{railwaymarket} Mordor Intelligence, Railway Management System Market -- Growth, Trends and Forecasts, Mordor Intelligence, 2024, \url{https://www.mordorintelligence.com/} (dostęp: 30.06.2026).

\bibitem{infrastructureascode} Morris K., Infrastructure as Code: Managing Servers in the Cloud, wyd. 2, O'Reilly Media, Sebastopol 2020.

\bibitem{cypher} Neo4j, Cypher Query Language Reference, \url{https://neo4j.com/docs/cypher-manual/} (dostęp: 30.06.2026).

\bibitem{openstreetmap} OpenStreetMap contributors, OpenStreetMap, \url{https://www.openstreetmap.org/} (dostęp: 30.06.2026).

\bibitem{pythonasyncio} Python Software Foundation, asyncio -- Asynchronous I/O. Python Documentation, \url{https://docs.python.org/3/library/asyncio.html} (dostęp: 30.06.2026).

\bibitem{fastapi} Ramírez S., FastAPI Documentation, \url{https://fastapi.tiangolo.com/} (dostęp: 30.06.2026).

\bibitem{graphdatabases} Robinson I., Webber J., Eifrem E., Graph Databases: New Opportunities for Connected Data, wyd. 2, O'Reilly Media, Sebastopol 2015.

\bibitem{haversine} Sinnott R.W., Virtues of the Haversine, Sky and Telescope, vol. 68, nr 2, 1984, str. 159.

\bibitem{sommerville2016} Sommerville I., Software Engineering, wyd. 10, Pearson, Boston 2016.

\bibitem{svelte} Svelte, Svelte Documentation, \url{https://svelte.dev/docs} (dostęp: 30.06.2026).

\bibitem{tailwind} Tailwind Labs, Tailwind CSS Documentation, \url{https://tailwindcss.com/docs} (dostęp: 30.06.2026).

\end{thebibliography}

\end{document}
